\subsection{Libratus Surprises the Poker Community with Overbets~\citep{BrownLibratus2017, KanjunAndNoam}}
\label{sec:libratus}

Multi-agent games with incomplete information~\citep{leytonbrown2008gametheory}, such as poker, present unique challenges for AI. Unlike games with perfect information like chess or Go, agents must contend with hidden variables such as an opponent's cards and psychological state, as well as the need for deception and bluffing. For years, these complexities were considered a significant barrier for AI, leading many to believe that human intuition would remain superior.

However, a team of researchers led by Dr. Noam Brown developed a neural-network-based poker agent known as \emph{Libratus} and trained it to compete at the highest levels of no-limit Texas Hold'em. Prior to Libratus, expert human players felt that AI techniques were still far from posing a genuine threat, particularly in games characterized by hidden information and bluffing. After Libratus defeated the top poker players in 2017, Brown recalled in an interview~\citep{KanjunAndNoam}:

\begin{quote}
The reception in the poker community was one of surprise and shock. People looked at the 2015 competition where [AI techniques] lost and thought ``oh, we are really far away from [these techniques] being a threat to poker'' and then to see these top expert players losing by the margin that they did [just one year later], I think really shocked the poker community. People were telling us they literally did not believe it was possible to beat top humans by the margin [Libratus was winning by]. And especially because of the way that [Libratus] played. At the end of each day, we gave the expert humans a log of all of the hands that the bot had in each hand of poker. That is golden information. If you are playing poker against somebody, you see maybe a third of the hands that actually reach showdown, and to just give somebody a log of all of the hands that the bot had and to still lose despite that information, a lot of people were really surprised by that.
I think people realized that there is a much higher [skill] ceiling [in poker] than they had previously thought.
\end{quote}

Reflecting on the long-term impact of Libratus and whether it prompted improvements in human play, Brown continues:

\begin{quote}
[P]rofessional poker players these days all use bots as training tools, it is a big industry now. And the game has changed a lot since the 2017 competition. One thing [Libratus] loved to do was these things called overbets; so humans, when they play poker, they size their bets relative to the size of the pot. So, if there is 100 dollars in the pot, maybe you will bet between 25 and 100 dollars and maybe if you are feeling really adventurous, you will bet 150 dollars. But the bot would sometimes bet around 10000 dollars. And the humans, when they saw this, it put them in really difficult positions. Like, [the humans] would have the second best hand that is possible sometimes, and then the bot just goes all-in on a 100 to 200 dollar pot and the human is just sitting there for five minutes thinking ``oh my god, do they have the best possible hand? Are they bluffing?'' And a key insight is if you see your opponents really struggling with a decision, you know you're playing good poker.
When I saw the experts really struggling with those kinds of decisions, I knew we were doing really well.

That technique was one of the main things that the experts walked away from the competition saying they would do more---this idea of applying overbets if it is done in the right way. A lot of bad players will bet all in also for no reason and that is not a good strategy which is why a lot of [top human] players ignored that strategy for a long time because it was associated with really bad play. But it turns out that if you are able to [overbet] in the right way, in the right spots, with the right balance of hands, it is an extremely effective strategy and has become much more popular in high-stakes poker these days [as a result of Libratus].
\end{quote}
